% !PRINTED 220
\weberpage{220}
%<<CONT

\weq{5}{y = t_1\Bigl(a_0x + \frac{a_1}{3}\Bigr) + t_0\Bigl(a_0x^2 + a_1x + \frac{2a_2}{3}\Bigr)}

für $y$ ergiebt. Schreiben wir die Gleichung in der Form

\[ y^3 + P_2y + P_3 = 0, \]

so sind $P_2, P_3$ Covarianten von $f$, und zwar $P_2$ von der zweiten, $P_3$ von der dritten Ordnung, sowohl in $t$ als in $a$.

Da nun $P_2$ und $P_3$ Covarianten der cubischen Form sind, so können sie sich nach \S.\,62 und 63 von den beiden dort definirten Functionen

\[ H(t_1,t_0), \qquad Q(t_1,t_0) \]

nur um constante, d.\,h. numerische Factoren unterscheiden. Diese constanten Factoren lassen sich durch irgend eine specielle Annahme bestimmen.

Wir können z.\,B. annehmen

\[ t_1 = 1, \quad t_0 = 0, \quad a_0 = 1, \quad a_1 = 0, \]

dann wird $y=x$, also

\[ P_2 = a_2, \qquad P_3 = a_3. \]

Andererseits ergiebt sich aber nach den im \S.\,62 gegebenen Formeln (2), (3), (7) für diese besondere Annahme

\[ H(1,0) = 3a_2, \qquad Q(1,0) = 27a_3, \]

woraus man allgemein schliesst

\[ P_2 = \frac{1}{3}H(t_1,t_0), \qquad P_3 = \frac{1}{27}Q(t_1,t_0), \]

so dass wir für $y$ die cubische Gleichung erhalten

\weq{6}{y^3 + \frac{1}{3}H(t_1,t_0)y + \frac{1}{27}Q(t_1,t_0) = 0.}

Die Discriminante $\Delta$ dieser cubischen Gleichung ist

\[ \Delta = -\frac{1}{27}(4H^3 + Q^2), \]

also mit Anwendung der Relation \S.\,62, (12)

\weq{7}{4H^3 + Q^2 + 27Df^2 = 0}

\weq{8}{\Delta = Df(t_1,t_0)^2,}

wenn $D$ die Discriminante der gegebenen cubischen Gleichung ist, in Uebereinstimmung mit den allgemeinen Resultaten des vorigen Paragraphen.

Wollen wir hierauf die Auflösung der cubischen Gleichung gründen, so müssen wir zunächst nach den Vorschriften des vorigen Paragraphen $x$ rational durch $y$ darstellen. Wir setzen

\weq{9}{a_0f(t_1,t_0)x = Q_0 + Q_1y + Q_2y^2.}
